% ============================================================
\section{Resultados e interpretação}\label{sec:resultados}
% ============================================================

Salvo indicação, os valores são da configuração principal: $k=10$, desempate por sorteio,
rede por união nas métricas estruturais e direcionada nas vizinhanças. As curvas nos 36 blocos
vêm da segunda execução, com a mesma amostra e os mesmos vetores.

% ------------------------------------------------------------
\subsection{As quatro redes}\label{sec:res-geral}

A Tabela~\ref{tab:metricas-globais} reúne as métricas estruturais da fase final, que repetem as
da entrega parcial com a média sobre as 20 sementes do desempate e acrescentam a reciprocidade,
a assimetria do grau de entrada, as comunidades e a rede do vocabulário.

\tabelaopcional{metricas_globais}
    {Métricas estruturais das redes ($k=10$, versão por união). Na rede lexical, média $\pm$
    desvio entre as 20 sementes do desempate; nas camadas contextuais, semente 0.}
    {tab:metricas-globais}

A clusterização alta das redes \cod{lex} e \cod{L01} e os 44 componentes da rede lexical são
efeito da construção: uma ocorrência de um tipo frequente tem os 10 vizinhos entre as outras
ocorrências do mesmo tipo, que também se escolhem entre si e formam grupos densos, às vezes
isolados. Nos blocos 18 e 36, as vizinhanças misturam tokens diferentes, os triângulos rareiam,
os grupos se fundem e surgem atalhos. O número de comunidades cai de 186 para 58 e 44, e a
modularidade de 0,95 para 0,77: as comunidades ficam maiores e menos separadas.

\figuraopcionallado{graus_camadas}
    {Distribuição do grau de entrada nas redes direcionadas: fração dos vértices com grau de
    entrada pelo menos $d$, em escala log-log. A linha vertical marca $k=10$, a média comum a
    todas as redes. Quanto mais longa a cauda, mais fortes os \emph{hubs}.}
    {fig:graus-camadas}

\paragraph{\emph{Hubs} e anisotropia.} O grau de entrada é estreito na rede lexical (máximo de
30, assimetria próxima de zero), porque o sorteio espalha os arcos entre cópias idênticas
(Figura~\ref{fig:graus-camadas}). A cauda cresce nos blocos 1 e 18 (assimetria 1,0 e 1,4) e
fica pesada no bloco 36: um vértice (\tok{\esp apresent}) recebe 247 arcos e a assimetria
chega a 4,9; os maiores \emph{hubs} são inícios e continuações de palavras e dígitos.
\emph{Hubs} são comuns em espaços de alta
dimensão~\cite{radovanovic2010}, e aqui crescem junto com a anisotropia: o cosseno médio entre
pares de vértices ao acaso sobe de 0,03 na camada lexical para 0,39 no bloco 18 e 0,58 no
bloco 36. A reciprocidade sobe de 0,49 para 0,61 no bloco 18 e cai para 0,44 no 36, porque
muitos vértices escolhem um \emph{hub} que não os escolhe de volta.

\paragraph{Uma ilha de primeiros tokens.} O segundo componente das redes \cod{L18} e \cod{L36}
é formado pelo primeiro token de vários parágrafos (quase sempre ``O'' ou ``A''). No bloco 18,
os vetores da posição 1 têm norma cerca de 50 vezes maior que a dos outros (mediana de 2.814
contra 55), as \emph{ativações massivas}~\cite{sun2024}, e todos os seus vizinhos são outros
vértices da posição 1. São 0,6\% dos vértices, e retirá-los quase não muda as métricas globais.

% ------------------------------------------------------------
\subsection{P1: quanto e de que forma a organização muda}\label{sec:res-p1}

\textbf{Resposta:} a organização quase não muda no bloco 1 e muda muito nos blocos 18 e 36, mas
a identidade do token não desaparece.

\paragraph{Troca de vizinhos.} O Jaccard médio entre a vizinhança lexical e a de cada camada é
0,35 no bloco 1, 0,16 no bloco 18 e 0,10 no bloco 36: no último bloco, uma ocorrência mantém em
média menos de 2 dos seus 10 vizinhos lexicais. Descontado o piso de ruído, a mudança é de 0,14,
0,34 e 0,40. A Figura~\ref{fig:p1-jaccard-faixas} mostra por que o desconto importa. Para tipos
com 10 ou mais ocorrências, o Jaccard entre \cod{lex} e \cod{L01} é igual ao piso (0,38 contra
0,39 na faixa de 10 a 49; 0,12 contra 0,12 na de 50): no bloco 1, os vizinhos continuam sendo as
outras ocorrências do mesmo token, e só muda \emph{quais} cópias foram sorteadas (o Jaccard por
composição de tipos, que ignora o sorteio, vale 0,98). Já os tipos raros mudam desde o bloco 1
(Jaccard de 0,29 e 0,50, contra pisos de 0,73 e 0,76): os seus vizinhos lexicais eram outros
tipos parecidos, e o bloco 1, que já vê o token anterior, os troca. Nos blocos 18 e 36, todas
as faixas ficam bem abaixo do piso.

\figuraopcional{p1_jaccard_faixas}
    {Jaccard entre a vizinhança de cada ocorrência na rede lexical e nos blocos 1, 18 e 36, por
    faixa de frequência do tipo (caixas: quartis; marcadores: médias). A linha horizontal é o
    piso de ruído da faixa: só a distância abaixo dela é reorganização.}
    {fig:p1-jaccard-faixas}

\paragraph{Dominância lexical.} A fração de vizinhos do mesmo tipo é 0,70 na rede lexical (o
teto), 0,69 no bloco 1, 0,43 no bloco 18 e 0,25 no bloco 36; relativa ao teto, cai de 1 para
0,99, 0,62 e 0,40. Mesmo no bloco 36, uma ocorrência de um tipo com 50 ocorrências tem em média
3,9 vizinhos do mesmo token, quando o acaso daria 0,04. A queda depende da categoria: no bloco
36, a dominância relativa é de 0,25 nas continuações de palavra, 0,29 nos números, 0,35 nas
palavras inteiras, 0,54 na pontuação e 0,79 nos inícios de palavra. Os inícios de palavra são
a exceção: a dominância deles até sobe entre os blocos 18 e 36 (de 0,73 para 0,79). Uma
explicação provável é que o bloco 36 codifica o próximo token (Seção~\ref{sec:res-lente}), e o
que vem depois de \tok{\esp apresent} é a continuação da palavra, que depende sobretudo do
próprio início; duas ocorrências dele preveem continuações parecidas e ficam próximas.

\textbf{Interpretação.} A H1 se confirma: a dominância cai, mas fica muito acima do acaso;
surgem \emph{hubs} no bloco 36; e, como mostra a P3, as comunidades se alinham menos ao token.
A mudança não é só local: ela troca vizinhos, encurta distâncias, funde componentes e cria
\emph{hubs}.

% ------------------------------------------------------------
\subsection{P2: em que camadas a mudança se concentra}\label{sec:res-p2}

\textbf{Resposta:} a mudança se espalha pelo modelo, com três momentos mais fortes (os
primeiros blocos, os blocos 22 a 25 e o último), e a ordem depende da categoria do token.

O painel da esquerda da Figura~\ref{fig:p2-curvas-camadas} mostra o Jaccard entre as
vizinhanças de blocos consecutivos: baixo no começo (0,54 entre os blocos 1 e 2), cerca de 0,83
nos blocos 15 a 18 (cada bloco mantém uns 9 dos 10 vizinhos), 0,67 no bloco 25 e 0,59 no último
bloco. Os blocos do meio mudam pouco de cada vez, mas a mudança acumulada é grande. O painel da
direita mostra a dominância lexical: as continuações de palavra caem de 0,61 para 0,26 já no
bloco 4, e os números de 0,93 para 0,60 no bloco 3; as palavras inteiras perdem a identidade
devagar (0,86 no bloco 1, 0,64 no bloco 20) e caem mais rápido depois do bloco 28 (0,30 no
bloco 36); os inícios de palavra ficam quase constantes. Entre as três transições da
configuração principal, a maioria dos vértices perde mais vizinhos entre os blocos 1 e 18
(48\%) ou entre 18 e 36 (37\%); as continuações, sobretudo na primeira (78\%), e as palavras
inteiras, na segunda (50\%).

\figuraopcional{p2_curvas_camadas}
    {Curvas nos 36 blocos. À esquerda, o Jaccard entre as vizinhanças de blocos consecutivos
    (no bloco 1, a comparação é com a rede lexical); à direita, a dominância lexical $D$ em cada
    bloco, por categoria de token.}
    {fig:p2-curvas-camadas}

\textbf{Interpretação.} A H2 se confirma em parte. Continuações de palavra e dígitos não
significam nada sozinhos e perdem a identidade nos primeiros blocos, quando o modelo junta os
pedaços; palavras inteiras mantêm a identidade até o último terço. Os inícios de palavra não
seguem a hipótese: quase não perdem a identidade.

% ------------------------------------------------------------
\subsection{P3: comunidades e contexto}\label{sec:res-p3}

\textbf{Resposta:} sim. As comunidades das redes \cod{lex} e \cod{L01} seguem o token; as do
meio do modelo seguem mais o texto de origem do que o token; no fim, o próximo token ganha peso.

Medida pelo NMI $-$ base, a concordância das comunidades com o token é 0,32 nas redes \cod{lex}
e \cod{L01}, e a com qualquer rótulo de contexto fica abaixo de 0,05. No bloco 18 a ordem se
inverte: artigo 0,25, parágrafo 0,24, sentença 0,21 e tema 0,20, contra 0,18 do token. No bloco
36, os rótulos de contexto caem para cerca da metade (tema 0,09, artigo 0,12), enquanto o
próximo token sobe de 0,08 para 0,19 e o token previsto pelo modelo para 0,22. A faixa de
posição nunca passa de 0,04: as comunidades não são um artefato da posição.

\figuratodas{p3_nmi_camadas}
    {Concordância (NMI $-$ base) entre as comunidades de cada bloco e cada rótulo, nos
    vértices do núcleo; o bloco 0 é a rede lexical. O último painel mostra a estabilidade do
    Leiden (concordância média entre as 10 execuções).}
    {fig:p3-nmi-camadas}

A Figura~\ref{fig:p3-nmi-camadas} mostra as mesmas medidas nos 36 blocos. A concordância com o
contexto acelera a partir do bloco 9, tem máximos nos blocos 12 e 13 (artigo 0,27, parágrafo
0,26, tema 0,22) e 19, cai bruscamente nos blocos 22 a 24 (o tema vai a 0,08), volta no 25 e
diminui até o 36. A concordância com o próximo token cresce a partir do bloco 21 e chega a 0,19
no bloco 34. A queda dos blocos 22 a 24 coincide com a do Jaccard entre blocos consecutivos e
com o bloco em que o próximo token passa a ficar mais perto do vetor do que o próprio token
(Seção~\ref{sec:res-lente}). A estabilidade do Leiden cai de 0,96 para 0,83: as comunidades das
últimas camadas são menos nítidas, mas reprodutíveis.

\textbf{Interpretação.} A H3 se confirma com uma correção: a organização pelo contexto não
cresce até o fim; ela tem um máximo nos blocos do meio, onde as comunidades reúnem ocorrências
do mesmo artigo e parágrafo independentemente do token, e cede lugar ao próximo token nos
últimos blocos. Isso concorda com a visualização da entrega parcial, em que as arestas entre
ocorrências do mesmo tema eram 52\% no bloco 18 e 44\% no 36.

% ------------------------------------------------------------
\subsection{P4: separação das ocorrências de um mesmo token}\label{sec:res-p4}

\textbf{Resposta:} sim. As ocorrências de um token, idênticas na camada lexical, se separam, e
as palavras ambíguas se separam mais do que as de controle, por tema e, mais ainda, pelo
sentido anotado.

\figuratodas{p4_separacao}
    {Separação das ocorrências de um mesmo tipo nos 36 blocos (médias por grupo, tipos com pelo
    menos 10 ocorrências). O painel do NMI mostra o valor bruto, sem descontar a linha de base,
    que difere entre os grupos (no bloco 36: 0,09 nos alvos, 0,13 nos controles, 0,26 nas
    multitema e 0,19 nas funcionais); o texto compara os grupos pelo NMI $-$ base.}
    {fig:p4-separacao}

\paragraph{Os grupos de palavras (Figura~\ref{fig:p4-separacao}).} A
\textbf{autossimilaridade} segue a mesma curva em todos os grupos: cai de 1 para 0,6 a 0,7 nos
blocos 10 a 15 e volta a cerca de 0,78 no bloco 36. Ela não distingue palavras ambíguas de
estáveis, e a subida final acompanha a anisotropia. A diferença \textbf{intra $-$ entre temas}
separa os grupos: nos alvos, ela tem o máximo nos blocos 11 e 12 (0,12) e vale 0,09 no bloco 18
e 0,05 no 36; nos controles, não passa de 0,02. O \textbf{NMI $-$ base entre comunidade e tema}
é de 0,23 nos alvos no bloco 18 e 0,33 no 36, contra 0,02 e 0,04 nos controles e nas palavras
multitema. No bloco 36, as ocorrências de um alvo se espalham por 3,3 comunidades efetivas, e as
de um controle por 2,7.

\tabelaopcional{p4_alvos}
    {P4 palavra por palavra, no bloco 36: separação pelo tema e pelo sentido anotado, em todas
    as ocorrências e só nas que o contexto à esquerda decide.}
    {tab:p4-alvos}

\paragraph{Palavra por palavra (Tabela~\ref{tab:p4-alvos}).} Todos os alvos têm diferença
intra $-$ entre temas positiva, de 0,024 (\emph{massa}) a 0,094 (\emph{órgão}); nos controles
ela fica perto de zero, com uma exceção, \emph{primeiro} (0,044): o tema sozinho, sem mudança de
sentido, também separa um pouco as ocorrências, e por isso os controles são necessários. Nas 11
palavras-alvo, a separação pelo sentido anotado é maior que pelo tema (\emph{família}: 0,118
pelo sentido e 0,027 pelo tema), e o NMI entre comunidade e sentido vai de 0,33 (\emph{banco})
a 0,70 (\emph{matriz}), contra cerca de 0,13 ao acaso. Nas ocorrências em que o texto à esquerda
já decide o sentido, os valores são parecidos ou maiores (\emph{nota} chega a 0,84): a
separação não depende de informação à direita, que o modelo não vê.

\paragraph{Onde os sentidos se separam.} No bloco 18, cinco palavras (\emph{órgão},
\emph{família}, \emph{capital}, \emph{matriz} e \emph{estado}) já dividem as ocorrências em
comunidades pelo sentido, e as outras seis ficam quase inteiras numa só comunidade. No bloco
36, todas se dividem, e as divisões seguem o sentido: a maior parte das ocorrências de
\emph{banco} de dados fica numa comunidade, e a das de \emph{banco} financeiro, em outras. A
diferença de cosseno entre sentidos é máxima nos blocos 11 e 12 (0,17 em média) e cai para
0,08 no bloco 36, enquanto o NMI entre comunidade e sentido sobe até o fim (0,37 no bloco 18 e
0,56 no 36, com linhas de base de 0,06 e 0,13).

\textbf{Interpretação.} A H4 se confirma: os alvos se separam mais que os controles em todas as
medidas que levam o tema ou o sentido em conta, e o sentido explica a separação melhor que o
tema. A geometria separa mais os sentidos no meio do modelo; a partição da rede segue mais o
sentido no fim, quando as ocorrências de cada palavra se espalham por mais comunidades.

% ------------------------------------------------------------
\subsection{Do próprio token ao próximo, e a rede do vocabulário}\label{sec:res-lente}

A Figura~\ref{fig:lente-ranking} mostra, bloco a bloco, a posição do próprio token e do
próximo token entre as linhas de $E$ ordenadas pela semelhança com o vetor da ocorrência. No
bloco 1, a linha mais próxima é a do próprio token em 90\% das ocorrências; no bloco 5, em
16\%; e a posição mediana dele chega a cerca de 30.000 entre os blocos 13 e 22. A posição
mediana do próximo token fica entre 50.000 e 70.000 até o bloco 21 e depois cai depressa
(2.100 no bloco 30, 5 no 36); no bloco 36, o próximo token é a linha mais próxima em 31\% das
ocorrências e está entre as 20 mais próximas em 69\%. A partir do bloco 23, ele fica mais perto
do vetor do que o próprio token na maioria das ocorrências. Pela matriz $E$, a identidade sai do
vetor bem antes do que a dominância lexical sugere: as ocorrências de um token continuam
vizinhas entre si (P1) mesmo quando já não estão perto da linha do token.

\figuraopcional{lente_ranking}
    {Posição mediana do próprio token (à esquerda) e do próximo token (à direita) entre as
    linhas da matriz de \emph{embeddings}, ordenadas pelo cosseno com o vetor da ocorrência,
    por bloco e categoria de token (escala logarítmica; 1 = a linha mais próxima).}
    {fig:lente-ranking}

A rede do vocabulário (última coluna da Tabela~\ref{tab:metricas-globais}) é conexa, mas tem
\emph{hubs} extremos (grau de entrada máximo de 1.756, assimetria de 22,6): caracteres isolados
de escritas raras, provavelmente pouco vistos no treinamento. As suas 171 comunidades quase não
seguem a escrita (NMI de 0,11) e misturam línguas: os vizinhos de cada palavra-alvo são
variantes de grafia e traduções para o inglês, o espanhol e o chinês, e indicam um só sentido
(os de \emph{banco}, o financeiro). Na entrada, uma palavra ambígua fica perto de um dos seus
sentidos; os outros só aparecem com o contexto.

% ------------------------------------------------------------
\subsection{Robustez}\label{sec:robustez}

Refizemos as redes mudando uma escolha de cada vez, e \textbf{em nenhuma variação a conclusão
qualitativa muda}: a clusterização é maior no bloco 1 que nos blocos 18 e 36, as distâncias
encurtam a partir do bloco 18 e a troca de vizinhos em relação à rede lexical cresce do bloco 1
ao 36. Com $k=5$ e $k=20$, os valores mudam na direção esperada (distância média no bloco 18 de
7,3 e 4,0, contra 5,1 com $k=10$), mas a ordem entre as camadas se mantém. O mesmo vale para a
rede mútua, para o desempate pela posição (Jaccard de 0,359 contra 0,354 entre \cod{lex} e
\cod{L01}), para a aceitação de todos os empatados e para uma tolerância cem vezes maior. A
saída normalizada do bloco 36 dá redes quase iguais às da saída bruta, e subtrair a média dos
vetores, o que remove a direção comum, muda pouco: a anisotropia não produz os resultados. Com
outras resoluções do Leiden e com o Louvain, a modularidade de \cod{L18} e \cod{L36} fica
entre 0,73 e 0,76. A P3 e a P4 não foram repetidas nessas variações.
